% Copyright 2026 MarcosHCK
%
\documentclass[12pt,a4paper]{article}
\usepackage{amsfonts}
\usepackage{amsmath}
\usepackage{amsthm}
\usepackage{changepage}
\usepackage{hyperref}
\usepackage{bookmark} % must load after hyperref
\usepackage{cleveref} % must load after hyperref
\usepackage{float}
\usepackage{graphicx}
\usepackage[utf8]{inputenc}
\usepackage{indentfirst}
\usepackage{polyglossia} % languages set below via \setmainlanguage/\setotherlanguage
\usepackage{setspace}
\usepackage[dvipsnames]{xcolor}

\setmainlanguage{spanish}
\setotherlanguage{english}

\setstretch{1.5}

\def\mathdefault#1{#1}
\newcommand{\myeoch}{\pagebreak}

\providecolor{Aquamarine}{rgb}{0.498, 1, 0.83}
\providecolor{DarkAquamarine}{rgb}{0.18, 0.55, 0.34}
\providecolor{YellowOrange}{rgb}{1, 0.7, 0}
\newcommand\myshade{85}
\colorlet{mycitecolor}{YellowOrange}
\colorlet{mylinkcolor}{violet}
\colorlet{myurlcolor}{DarkAquamarine}

\hypersetup {
    colorlinks = true,
    citecolor  = mycitecolor!\myshade!black,
    linkcolor  = mylinkcolor!\myshade!black,
    urlcolor   = myurlcolor!\myshade!black,
  }

\begin{document}

\begin{titlepage}

  \centering
  {\huge\bfseries Reporte \#6 --- Parte 1\par}
	{\Large Autor: \itshape Marcos Antonio Pérez Lorenzo\par}
\end{titlepage}

\pagebreak

\section{Perceptrón sigmoide para una votación ponderada}

\subsection{Descripción de la red}

El archivo \texttt{perceptron.m} implementa una red neuronal de un perceptrón. Sus siete entradas representan siete votos binarios (1 para sí, 0 para no) con pesos \([1,2,1,2,1,2,1]\). La expresión matemática es:
\begin{equation*}
  y = \sigma\!\left(\sum_{i=1}^{7} w_i x_i + b\right)
    = \frac{1}{1 + e^{-\left(\sum_{i=1}^{7} w_i x_i + b\right)}}.
\end{equation*}
La salida es un valor entre cero y uno. Para tomar la decisión final, se utiliza el valor umbral \(0.5\): una salida igual o superior se interpreta como sí y una inferior como no.

\subsection{Organización del programa}

El programa comienza definiendo la cantidad de votos, los pesos de referencia y el umbral. A continuación genera las \(2^7=128\) combinaciones posibles de votos y asigna a cada una su respuesta objetivo aplicando la regla ponderada. De esta manera, los ejemplos de entrenamiento cubren todas las entradas binarias posibles para esta tarea.

La siguiente parte inicializa los siete pesos y el sesgo en cero, y establece la tasa de aprendizaje y el límite de épocas. Durante el entrenamiento, calcula la suma ponderada y la salida sigmoide para cada ejemplo, y ajusta los parámetros. Al terminar cada época, vuelve a evaluar las 128 combinaciones para comprobar si todas quedan clasificadas correctamente.

\subsection{Entrenamiento}

El aprendizaje es supervisado: para cada combinación de votos ya se conoce la respuesta que dicta la regla de referencia. La neurona procesa los ejemplos uno por uno y actualiza los pesos y el sesgo mediante descenso por gradiente sobre el error cuadrático, usando la derivada de la sigmoide para calcular la corrección. La tasa de aprendizaje es \(0.5\), y se permiten hasta \(5000\) épocas. El entrenamiento se detiene antes si, al aplicar un umbral de \(0.5\), todas las combinaciones tienen la clase correcta.

\subsection{Resultados}

La ejecución terminó tras \(23\) de las \(5000\) épocas permitidas. Los pesos obtenidos fueron aproximadamente:
$$[1.263010,\,2.922768,\,1.396897,\,2.985553,\,1.421800,\,3.025049,\,1.428059]$$
Y sesgo
$$-7.466145$$

No es necesario que estos valores coincidan exactamente con los pesos de referencia, lo relevante es que reproduzcan las decisiones de la regla.

El programa clasificó correctamente las \(128\) combinaciones, es decir, obtuvo una exactitud de entrenamiento del \(100.00\%\). Entre los ejemplos impresos, la entrada \([0,0,0,0,0,0,0]\) produjo una salida sigmoide de \(0.000572\) y la decisión ``No''; la entrada \([1,1,1,1,1,1,1]\) produjo \(0.999068\) y ``Sí''. Para la entrada de demostración \([1,0,1,1,0,1,0]\), la suma ponderada aprendida fue \(8.670509\), la salida fue \(0.769300\) (decisión: sí).

\end{document}
